\label{sec:method-region}

\begin{figure}[tb]
\centering
\begin{tikzpicture}[
  stage/.style={draw, thick, rounded corners, minimum height=1.1cm, minimum width=2.6cm, fill=none},
  arrow/.style={-{Stealth[length=2.2mm]}, thick}
]
\node[stage, label={above:batch $b$ ($B$ ops)}] (batch) {};
\node[stage, right=0.55cm of batch, label={above:witness counts}] (wit) {};
\node[stage, right=0.55cm of wit, label={above:seed set $S$}] (seed) {};
\node[stage, right=0.55cm of seed, label={above:mate closure $R$}] (clo) {};
\node[stage, right=0.55cm of clo, label={above:region peel}] (peel) {};
\node[stage, below=0.75cm of peel, label={below:exact $\tau$}] (out) {};
\draw[arrow] (batch) -- node[above]{\small enumer\-ate triangles} (wit);
\draw[arrow] (wit) -- node[above]{\small sup-changed mates} (seed);
\draw[arrow] (seed) -- node[above]{\small parallel BFS} (clo);
\draw[arrow] (clo) -- node[above]{\small recount $+$ virtual pops} (peel);
\draw[arrow] (peel) -- (out);
\end{tikzpicture}
\caption{Processing one batch. All phases before the peel are parallel; the peel is one bounded sweep over the region.}
\label{fig:framework}
\end{figure}

\Cref{fig:framework} shows the three stages our algorithm applies to each batch: witness maintenance, mate closure, and a single region peeling with virtual boundary pops.
This section develops the first two stages and states the coverage guarantee; \Cref{sec:method-peel} develops the peeling stage and its correctness.

\paragraph{Witness maintenance}
Let $D$ be the deletions of the batch and $I$ its insertions, $b=D\uplus I$.
The witness phase updates the maintained supports $\sup(e)$ so that, after the phase, $\sup(e)$ equals the number of triangles containing $e$ in the \emph{current} graph, for every edge $e$ whose triangle set was touched by~$b$.
For every edge $f\in D\cup I$ we enumerate the common neighbors of its endpoints in parallel (each batch edge is an independent task), obtaining the set $\Delta(f)$ of triangles that are created or destroyed by the update.
For $f\in I$ we increment $\sup(g)$ by one for each $g$ in each enumerated triangle $\{f,g,h\}$; for $f\in D$ we decrement.
The per-mate contributions are aggregated with fetch-and-add on a per-batch counter array $\mathrm{dcount}$, so two batch updates that touch the same mate aggregate into a single delta rather than two traversals.
Deletions are processed before insertions, which makes the batch semantics unambiguous when a batch both deletes and re-inserts an edge: the deletion path removes the edge's id from the adjacency of its endpoints and the insertion path allocates a fresh id.
Because the triangle enumeration of phase 1 is the only place where the batch touches the global graph structure, its cost is $O\!\big(\sum_{f\in b}(\deg(f_1)+\deg(f_2))\big)$ work spread over $T$ threads.

\paragraph{Precise seeding}
After the witness phase, the algorithm knows every edge whose support changed: the batch edges themselves, and every mate whose $\mathrm{dcount}$ is nonzero.
These are the \emph{seeds} $S$.
A deleted edge may be seeded as a mate of another batch edge, or be a seed through its own deletion; only edges still present in the graph are seeded, and deleted edges are collected into the seed scan before their removal from the adjacency, so a mate that is itself deleted is dropped by the alive-guard at seeding time.
An obvious alternative---seeding every edge enumerated by phase 1, including mates whose support did not change---enlarges the seed set without enlarging the guaranteed change set.
Our implementation supports both modes (\texttt{--ablate allseeds}) and the experiments (\Cref{sec:experiments}) show that the analysis is tight: the two modes produce identical regions on every graph we measure, so the per-edge support check costs nothing in region quality.

\paragraph{Region closure}
The region is a superset of the truss-connected closure of $S$.
Starting from $S$, the closure performs a parallel breadth-first expansion over triangles: an unmarked edge $f$ is enqueued when it is a truss mate of a region edge $g$ on some triangle $\{g,f,h\}$, in which case $f$ and $h$ both join the region.
The expansion terminates because each edge enters the region at most once.
We prove below that this closure contains every edge whose trussness can change, so restricting all subsequent work to it is sound.

\begin{lemma}[Coverage]
\label{lem:coverage}
Let $G'$ be the graph obtained from $G$ by applying batch $b$, and let $C$ be the truss-connected mate closure of the seed set $S$ in $G'$.
If $\tau_{G'}(e)\ne \tau_G(e)$ for some edge $e$, then $e\in C$.
\end{lemma}

\begin{proof}
Run the peeling of \Cref{def:peeling} on $G$ and on $G'$, and let $e$ be an edge whose values differ.
Consider the first time the two peelings assign $e$ different values; without loss of generality the pops differ, so some mate $g$ of $e$ is popped at a different sweep position in the two runs.
Following the pop positions backwards, there is a first edge whose pop position differs because its own support sequence differs; supports differ only when a triangle witness is created or destroyed (a batch edge, hence in $S$) or when a mate with an earlier position difference decrements it.
Tracing this dependency backwards through mate adjacency yields a chain of edges from $e$ to a seed, in which consecutive edges share a triangle---that is, a mate-adjacency path from $e$ into $S$, so $e\in C$.
\end{proof}

Two structural remarks are in order.
First, the closure may be much smaller than the whole graph but much larger than the batch: in graphs with a single large truss-connected dense component every seed's closure reaches the whole component, and our algorithm then degenerates gracefully to peeling that component once per batch---still one traversal where per-edge algorithms perform $B$.
Second, the closure is computed in $G'$ (the post-batch graph), which keeps the expansion's adjacency scans simple: every edge visible during the expansion exists in the current graph.
